% modeling/part_intro-DE.tex
\chapter*{Modellierung}
\phantomsection
\addcontentsline{toc}{chapter}{Modellierung}

Dieser Teil legt dar, wie kanonische Eisenbahnsemantik als kompakte,
rekonstruierbare und operativ nutzbare Informationsstrukturen kodiert wird.

\section*{Was wissen wir bereits?}

Die vorangehenden Teile haben die begrifflichen Grundlagen von AIM geschaffen.
Intrinsische Geometrie, Eisenbahnzustände, Laufzeitbewertung, Realisierung und
realitätsbewusste Interpretation wurden als zusammenhängende theoretische
Konstrukte entwickelt.

Diese Begriffe definieren, was Eisenbahn-Alignments sind, wie sie sich
entwickeln und wie sie mit der physischen Welt zusammenwirken. Begriffe allein
definieren jedoch noch nicht, wie Alignment-Information in einem operativen
Informationsmodell dargestellt wird.

\section*{Was fehlt noch?}

Ein Informationsmodell erfordert ausdrückliche interne Strukturen.
Geometrische Begriffe müssen in Darstellungen überführt werden, die gespeichert,
ausgetauscht, rekonstruiert, analysiert und optimiert werden können.

Traditionelle Eisenbahnsysteme verwenden häufig Koordinatenlisten,
geometrische Primitive oder implementierungsspezifische Datenstrukturen. Für
Speicherzwecke mögen solche Darstellungen genügen; häufig verdecken sie jedoch
die intrinsische Struktur des Alignments und begrenzen die nachfolgende
Analyse.

Es fehlt daher ein Modellierungsrahmen, der Eisenbahnsemantik erfasst und
zugleich kompakt, interpretierbar und rechnerisch effizient bleibt.

\section*{Was folgt?}

Dieser Teil führt die Modellierungsebene von AIM ein. Zentrales Ziel ist die
Konstruktion dünn besetzter, semantischer Alignment-Modelle.

In AIM wird ein Alignment nicht primär durch abgetastete Koordinaten
dargestellt. Es wird vielmehr durch eine kompakte Sammlung geometrischer
Zustände, Übergangsbeschreibungen und semantischer Beziehungen dargestellt,
aus denen die Geometrie bei Bedarf rekonstruiert werden kann.

Die folgenden Kapitel führen ein:
\begin{itemize}
\item dünn besetzte Alignment-Darstellungen,
\item Übergangsmodellierung,
\item Übergangsklassifikation,
\item Alignment-Semantik höherer Ebene und
\item schwerpunktbasierte Modellierungsbegriffe.
\end{itemize}

Diese Strukturen bilden den operativen Kern des Alignment-Based Information
Modelling.

Sie zeigen zugleich, weshalb das System über die ursprüngliche Mathematik
hinauswuchs. Berlinish stellt eine verallgemeinerte Übergangsgrammatik bereit.
axtranNew verwendet Funktionen, Randbedingungen und Freiheitsgrade, um
konkrete Alignment-Strukturen zu berechnen. transitionDB bewahrt
wiederverwendbares Übergangswissen, und TransEd macht dieses Wissen
untersuchbar. SPOT wird anschließend benötigt, um Identität und Zustand
konkreter Ingenieurobjekte zu tragen, während alignmentOS die operative
Umgebung bereitstellt, in der die berechnete Struktur bearbeitet, bewertet und
mit anderer Arbeit verbunden werden kann. Dies sind voneinander abhängige
Antworten auf ein Ingenieurproblem und kein flacher Produktkatalog.

\section*{Bezug zur AIM-Roadmap}

In der AIM-Roadmap entspricht dieser Teil dem Übergang von begrifflichen
Eisenbahnbeschreibungen zu ausdrücklichen Alignment-Modellen.

\begin{figure}[ht]
\centering
% figures/modeling/icon_modeling_constructive_sparse-DE.tex
\begin{tikzpicture}[
  node distance=1.3cm,
  box/.style={draw, rounded corners, minimum width=4.3cm, minimum height=0.9cm, align=center},
  arrow/.style={->, thick}
]
\node[box] (elements) {Konstruktive Elemente};
\node[box, below=of elements] (sparse) {Dünn besetztes Alignment};
\node[box, below=of sparse] (model) {Engineering-Modell};
\draw[arrow] (elements) -- (sparse);
\draw[arrow] (sparse) -- (model);
\draw[thin] ($(sparse.center)+(-2.55,0.0)$) rectangle ($(sparse.center)+(2.55,2.6)$);
\end{tikzpicture}

\caption{Modellierungssymbol: Konstruktive Elemente bilden dünn besetzte Alignment- und Engineering-Modellsemantik.}
\label{fig:icon-modeling-constructive-sparse}
\end{figure}

Die Abfolge
\[
\text{Begriffe}
\rightarrow
\text{dünn besetztes Modell}
\rightarrow
\text{Alignment-Modell}
\]
markiert den Übergang von theoretischen Eisenbahnbeschreibungen zu operativen
Informationsstrukturen.

Die Abbildung beantwortet die Orientierungsfrage dieses Übergangs: Welche
konstruktiven Elemente müssen ausdrücklich erhalten bleiben, damit dünn
besetzte Modelle zugleich kompakt und semantisch treu bleiben?

Diese Modelle bilden anschließend die Grundlage für Laufzeitdienste,
Analysepipelines, Optimierungsverfahren und Engineering-Anwendungen.

\begin{principle}[Prinzip dünn besetzter Modellierung]
\label{princ:sparse-modeling}
In AIM sollen Eisenbahn-Alignments durch kompakte semantische Strukturen
dargestellt werden, aus denen geometrische, operative und
realisierungsbewusste Information bei Bedarf rekonstruiert werden kann.
\end{principle}

\begin{summary}
Begriffe erklären die Eisenbahnstruktur.

Realität erklärt das Zusammenwirken mit der physischen Welt.

Modellierung erklärt, wie Eisenbahnwissen dargestellt wird.

Dieser Teil begründet damit die dünn besetzten und semantischen
Informationsstrukturen, welche den Kern des Alignment-Based Information
Modelling bilden.
\end{summary}

Mit dieser Modellierungsebene kann der Teil Analyse untersuchen, wie diese
Strukturen durch Pipelines, Dienste und Systeminteraktionen transformiert
werden.
